\hypertarget{voynich-phytoglyphica-engine-specification}{%
\section{Voynich Phytoglyphica --- Engine
Specification}\label{voynich-phytoglyphica-engine-specification}}

\hypertarget{purpose}{%
\subsection{Purpose}\label{purpose}}

The Voynich Manuscript (VMS) has resisted decipherment for a century
because it is not a natural language text. It is a pharmaceutical
instruction set --- an executable split across six register sections,
where the opcode sequence lives in the recipe section (f103r+) and the
parameter values live nowhere in the manuscript at all. They were
supplied by the practitioner, using a structural grammar that the
manuscript presupposes but never states.

The Voynich Phytoglyphica project reconstructs that grammar. Using the
Imscribing Grammar (IG), every botanical entry in the VMS pharmacy
catalog can be assigned a 12-primitive structural tuple. That tuple
supplies the missing parameters: what solvent, what ratio, what
temperature, what comminution, how many cycles, when to stop. The VMS
opcodes tell you what operations to perform. The IG tuple tells you
exactly how.

The result is a fully elaborated pharmaceutical protocol for any VMS
botanical entry, derived entirely from structural analysis --- no
linguistic decipherment required.

\begin{center}\rule{0.5\linewidth}{0.5pt}\end{center}

\hypertarget{the-vms-as-split-register-executable}{%
\subsection{The VMS as Split-Register
Executable}\label{the-vms-as-split-register-executable}}

The manuscript is organized into six section-registers, each holding a
distinct component of the computation:

\begin{longtable}[]{@{}lll@{}}
\toprule
Section & Folios & Register Function\tabularnewline
\midrule
\endhead
Botanical & f1--f66 & Plant catalog --- identity index\tabularnewline
Pharmaceutical & f99--f102 & Pharmacy address --- operation
class\tabularnewline
Balneological & f75--f84 & Containment heap --- vessel
validation\tabularnewline
Astronomical & f67--f73 & Winding register --- cycle
authority\tabularnewline
Cosmological & f68 (foldout) & Invariant initialization --- Φ and
Ħ\tabularnewline
Recipe & f103r+ & Opcode sequence --- instruction stream\tabularnewline
\bottomrule
\end{longtable}

A complete computation requires all six sections. The session engine
reads them in order: cosmological initialization → pharmacy address →
heap validation → winding verification → recipe output → protocol
elaboration.

No section alone is sufficient. The recipe section contains only
abstract opcodes. The pharmacy section contains only address pointers.
The botanical catalog contains only identity entries. The grammar
supplies the parameter values that bind them.

\begin{center}\rule{0.5\linewidth}{0.5pt}\end{center}

\hypertarget{the-imscribing-grammar-12-primitives}{%
\subsection{The Imscribing Grammar --- 12
Primitives}\label{the-imscribing-grammar-12-primitives}}

The IG assigns every object in its catalog a 12-primitive structural
tuple. For botanical entries, each primitive encodes a distinct
pharmaceutical parameter. The 12 primitives, their canonical names, and
their roles in the pharmaceutical context:

\begin{longtable}[]{@{}lll@{}}
\toprule
Primitive & Name & Pharmaceutical Role\tabularnewline
\midrule
\endhead
Ð & Dimensionality & Registration depth --- structural address
class\tabularnewline
Þ & Topology & Plant material specification\tabularnewline
Ř & Recognition & Pattern-completion class --- extraction
completeness\tabularnewline
Φ & Parity & Solvent system and polarity\tabularnewline
ƒ & Fidelity & Yield and concentration target\tabularnewline
Ç & Kinetics & Extraction process, temperature, duration\tabularnewline
Γ & Granularity & Comminution specification\tabularnewline
ɢ & Coupling & Fraction combination mode\tabularnewline
⊙ & Criticality & Endpoint criterion\tabularnewline
Ħ & Chirality & Clarification and separation protocol\tabularnewline
Σ & Stoichiometry & Drug-to-solvent mass ratio\tabularnewline
Ω & Winding & Number of extraction cycles\tabularnewline
\bottomrule
\end{longtable}

Tuple notation: ⟨Ð⋅Þ⋅Ř⋅Φ⋅ƒ⋅Ç⋅Γ⋅ɢ⋅⊙⋅Ħ⋅Σ⋅Ω⟩. Values are Shavian characters
from the 49-symbol set (Shavian alphabet + ⊙). The primitive glyphs (Þ,
Ð, Φ, etc.) name the primitive as an entity; the Shavian value at each
position is its structural type in a given catalog entry.

\begin{center}\rule{0.5\linewidth}{0.5pt}\end{center}

\hypertarget{primitive-to-protocol-mapping}{%
\subsection{Primitive-to-Protocol
Mapping}\label{primitive-to-protocol-mapping}}

\hypertarget{uxfe-topology-plant-material}{%
\subsubsection{Þ --- Topology → Plant
Material}\label{uxfe-topology-plant-material}}

The Topology primitive encodes the structural organization of the
botanical object. In pharmaceutical terms: which part of the plant, and
how it is characterized.

\begin{longtable}[]{@{}lll@{}}
\toprule
\begin{minipage}[b]{0.30\columnwidth}\raggedright
Shavian\strut
\end{minipage} & \begin{minipage}[b]{0.30\columnwidth}\raggedright
Material\strut
\end{minipage} & \begin{minipage}[b]{0.30\columnwidth}\raggedright
Description\strut
\end{minipage}\tabularnewline
\midrule
\endhead
\begin{minipage}[t]{0.30\columnwidth}\raggedright
𐑡\strut
\end{minipage} & \begin{minipage}[t]{0.30\columnwidth}\raggedright
aerial parts\strut
\end{minipage} & \begin{minipage}[t]{0.30\columnwidth}\raggedright
leaf and stem, freshly dried\strut
\end{minipage}\tabularnewline
\begin{minipage}[t]{0.30\columnwidth}\raggedright
𐑰\strut
\end{minipage} & \begin{minipage}[t]{0.30\columnwidth}\raggedright
root / rhizome\strut
\end{minipage} & \begin{minipage}[t]{0.30\columnwidth}\raggedright
underground organs, cleaned and sliced\strut
\end{minipage}\tabularnewline
\begin{minipage}[t]{0.30\columnwidth}\raggedright
𐑥\strut
\end{minipage} & \begin{minipage}[t]{0.30\columnwidth}\raggedright
whole plant\strut
\end{minipage} & \begin{minipage}[t]{0.30\columnwidth}\raggedright
including seed heads and root crown\strut
\end{minipage}\tabularnewline
\begin{minipage}[t]{0.30\columnwidth}\raggedright
𐑶\strut
\end{minipage} & \begin{minipage}[t]{0.30\columnwidth}\raggedright
bark / pericarp\strut
\end{minipage} & \begin{minipage}[t]{0.30\columnwidth}\raggedright
outer cortex or fruit rind, dried\strut
\end{minipage}\tabularnewline
\begin{minipage}[t]{0.30\columnwidth}\raggedright
𐑸\strut
\end{minipage} & \begin{minipage}[t]{0.30\columnwidth}\raggedright
flowering tops\strut
\end{minipage} & \begin{minipage}[t]{0.30\columnwidth}\raggedright
whole entity in full anthesis; holographic self-similarity
preserved\strut
\end{minipage}\tabularnewline
\bottomrule
\end{longtable}

\hypertarget{ux3c6-parity-solvent-system}{%
\subsubsection{Φ --- Parity → Solvent
System}\label{ux3c6-parity-solvent-system}}

Parity encodes the bilateral balance between aqueous and non-aqueous
phases. The solvent system follows directly from the symmetry class of
the parity value.

\begin{longtable}[]{@{}lll@{}}
\toprule
Shavian & Solvent & Description\tabularnewline
\midrule
\endhead
𐑗 & water & 100\% aqueous, pH 6--7\tabularnewline
𐑿 & dilute aqueous & 5--10\% ethanol v/v in water\tabularnewline
𐑬 & hydroethanolic & 45--55\% ethanol v/v in water (bilateral
bridge)\tabularnewline
𐑯 & anhydrous ethanol & \textgreater95\% ethanol v/v\tabularnewline
𐑹 & fixed oil / CO₂ & cold-pressed carrier oil or supercritical CO₂
extract\tabularnewline
\bottomrule
\end{longtable}

Φ=𐑬 is the bilateral solvent: it bridges aqueous and lipophilic phases
simultaneously, which is why it is the most common value in the
nine-entry inaugural set. Φ=𐑹 (oil/CO₂) is the most non-polar,
appropriate for lipophilic constituents and saturated extraction
protocols.

\hypertarget{uxe7-kinetics-extraction-process}{%
\subsubsection{Ç --- Kinetics → Extraction
Process}\label{uxe7-kinetics-extraction-process}}

Kinetics encodes the thermodynamic rate-regime of the extraction.
Temperature and duration follow from the kinetic class.

\begin{longtable}[]{@{}lll@{}}
\toprule
\begin{minipage}[b]{0.30\columnwidth}\raggedright
Shavian\strut
\end{minipage} & \begin{minipage}[b]{0.30\columnwidth}\raggedright
Process\strut
\end{minipage} & \begin{minipage}[b]{0.30\columnwidth}\raggedright
Description\strut
\end{minipage}\tabularnewline
\midrule
\endhead
\begin{minipage}[t]{0.30\columnwidth}\raggedright
𐑘\strut
\end{minipage} & \begin{minipage}[t]{0.30\columnwidth}\raggedright
infusion\strut
\end{minipage} & \begin{minipage}[t]{0.30\columnwidth}\raggedright
single-pass ambient, 5--10 min, 20--25 °C\strut
\end{minipage}\tabularnewline
\begin{minipage}[t]{0.30\columnwidth}\raggedright
𐑤\strut
\end{minipage} & \begin{minipage}[t]{0.30\columnwidth}\raggedright
cold maceration\strut
\end{minipage} & \begin{minipage}[t]{0.30\columnwidth}\raggedright
12--24 h at 15--20 °C; frozen-order kinetics --- no heat\strut
\end{minipage}\tabularnewline
\begin{minipage}[t]{0.30\columnwidth}\raggedright
𐑧\strut
\end{minipage} & \begin{minipage}[t]{0.30\columnwidth}\raggedright
decoction\strut
\end{minipage} & \begin{minipage}[t]{0.30\columnwidth}\raggedright
15--30 min at 85--95 °C; activated kinetics --- sustained heat\strut
\end{minipage}\tabularnewline
\begin{minipage}[t]{0.30\columnwidth}\raggedright
𐑪\strut
\end{minipage} & \begin{minipage}[t]{0.30\columnwidth}\raggedright
percolation\strut
\end{minipage} & \begin{minipage}[t]{0.30\columnwidth}\raggedright
slow gravity-driven percolation at ambient\strut
\end{minipage}\tabularnewline
\begin{minipage}[t]{0.30\columnwidth}\raggedright
𐑺\strut
\end{minipage} & \begin{minipage}[t]{0.30\columnwidth}\raggedright
distillation\strut
\end{minipage} & \begin{minipage}[t]{0.30\columnwidth}\raggedright
steam or vacuum distillation; phase-separated barrier kinetics\strut
\end{minipage}\tabularnewline
\bottomrule
\end{longtable}

\textbf{Critical constraint (Ç=𐑤):} Cold maceration entries must never
be heated. When a VMS recipe step calls for Calefac (heating), that step
applies to an adjunct ingredient or excipient --- the carrier, the base,
the binding medium --- never to the primary cold extract. This is
enforced at the annotation level by detecting the frozen-order kinetics
flag.

\hypertarget{ux3c3-stoichiometry-drugsolvent-ratio}{%
\subsubsection{Σ --- Stoichiometry → Drug:Solvent
Ratio}\label{ux3c3-stoichiometry-drugsolvent-ratio}}

\begin{longtable}[]{@{}lll@{}}
\toprule
Shavian & Ratio & Description\tabularnewline
\midrule
\endhead
𐑙 & 1 : 1 & 1 g plant material per 1 mL solvent (saturated
loading)\tabularnewline
𐑕 & 1 : 2 & 1 g plant material per 2 mL solvent\tabularnewline
𐑳 & 1 : 3 & 1 g plant material per 3 mL solvent (triadic
ratio)\tabularnewline
\bottomrule
\end{longtable}

These are mass-to-volume ratios for the drug charge per extraction
cycle. At Σ=𐑙, the solvent is loaded to saturation --- the maximum
mass-transfer gradient is sustained throughout the extraction.

\hypertarget{ux3b3-granularity-comminution}{%
\subsubsection{Γ --- Granularity →
Comminution}\label{ux3b3-granularity-comminution}}

Granularity encodes the surface exposure class of the comminuted
material. Finer powder increases surface area and extraction efficiency
but can also extract unwanted components.

\begin{longtable}[]{@{}lll@{}}
\toprule
Shavian & Mesh & Description\tabularnewline
\midrule
\endhead
𐑚 & coarse & 2--4 mm pieces; no further comminution\tabularnewline
𐑔 & medium & pass mesh 40 (355 μm); coarser residue
discarded\tabularnewline
𐑲 & fine & pass mesh 100 (150 μm); uniform surface
exposure\tabularnewline
\bottomrule
\end{longtable}

\hypertarget{ux3c9-winding-extraction-cycles}{%
\subsubsection{Ω --- Winding → Extraction
Cycles}\label{ux3c9-winding-extraction-cycles}}

Winding encodes the cyclic structure of the extraction. The number of
cycles determines extraction depth and corresponds to the algebraic
period of the winding.

\begin{longtable}[]{@{}lll@{}}
\toprule
Shavian & Cycles & Mathematical class\tabularnewline
\midrule
\endhead
𐑷 & 1 & trivial winding (single pass)\tabularnewline
𐑴 & 2 & binary winding (ℤ₂ period)\tabularnewline
𐑭 & 3 & integer winding (ℤ period; three complete turns)\tabularnewline
𐑟 & continuous & non-Abelian winding (braid-group period; percolation
class)\tabularnewline
\bottomrule
\end{longtable}

Each cycle: fresh solvent charge on the marc from the previous cycle.
Combined fractions at Compone/endpoint. The winding count governs total
solute recovery.

\hypertarget{ux192-fidelity-concentration-target}{%
\subsubsection{ƒ --- Fidelity → Concentration
Target}\label{ux192-fidelity-concentration-target}}

Fidelity encodes the yield relationship between raw extract and final
preparation. Higher fidelity requires volume reduction and marker
standardization.

\begin{longtable}[]{@{}lll@{}}
\toprule
\begin{minipage}[b]{0.30\columnwidth}\raggedright
Shavian\strut
\end{minipage} & \begin{minipage}[b]{0.30\columnwidth}\raggedright
Target\strut
\end{minipage} & \begin{minipage}[b]{0.30\columnwidth}\raggedright
Description\strut
\end{minipage}\tabularnewline
\midrule
\endhead
\begin{minipage}[t]{0.30\columnwidth}\raggedright
𐑱\strut
\end{minipage} & \begin{minipage}[t]{0.30\columnwidth}\raggedright
1× (standard)\strut
\end{minipage} & \begin{minipage}[t]{0.30\columnwidth}\raggedright
no reduction; proportional yield; linear fidelity\strut
\end{minipage}\tabularnewline
\begin{minipage}[t]{0.30\columnwidth}\raggedright
𐑞\strut
\end{minipage} & \begin{minipage}[t]{0.30\columnwidth}\raggedright
2× (concentrated)\strut
\end{minipage} & \begin{minipage}[t]{0.30\columnwidth}\raggedright
reduce to half volume after extraction; quadratic fidelity\strut
\end{minipage}\tabularnewline
\begin{minipage}[t]{0.30\columnwidth}\raggedright
𐑐\strut
\end{minipage} & \begin{minipage}[t]{0.30\columnwidth}\raggedright
3× (highly concentrated)\strut
\end{minipage} & \begin{minipage}[t]{0.30\columnwidth}\raggedright
reduce to one-third volume; cubic fidelity; standardize to marker
compound\strut
\end{minipage}\tabularnewline
\bottomrule
\end{longtable}

\hypertarget{ux127-chirality-clarification-protocol}{%
\subsubsection{Ħ --- Chirality → Clarification
Protocol}\label{ux127-chirality-clarification-protocol}}

Chirality encodes the stereochemical resolution class of the
clarification step. Higher chirality demands separate the extract into
increasingly refined fractions.

\begin{longtable}[]{@{}lll@{}}
\toprule
\begin{minipage}[b]{0.30\columnwidth}\raggedright
Shavian\strut
\end{minipage} & \begin{minipage}[b]{0.30\columnwidth}\raggedright
Protocol\strut
\end{minipage} & \begin{minipage}[b]{0.30\columnwidth}\raggedright
Description\strut
\end{minipage}\tabularnewline
\midrule
\endhead
\begin{minipage}[t]{0.30\columnwidth}\raggedright
𐑓\strut
\end{minipage} & \begin{minipage}[t]{0.30\columnwidth}\raggedright
none\strut
\end{minipage} & \begin{minipage}[t]{0.30\columnwidth}\raggedright
use as-is; racemic; no chiral resolution\strut
\end{minipage}\tabularnewline
\begin{minipage}[t]{0.30\columnwidth}\raggedright
𐑒\strut
\end{minipage} & \begin{minipage}[t]{0.30\columnwidth}\raggedright
single-step\strut
\end{minipage} & \begin{minipage}[t]{0.30\columnwidth}\raggedright
filter through coarse cloth or paper\strut
\end{minipage}\tabularnewline
\begin{minipage}[t]{0.30\columnwidth}\raggedright
𐑖\strut
\end{minipage} & \begin{minipage}[t]{0.30\columnwidth}\raggedright
two-step\strut
\end{minipage} & \begin{minipage}[t]{0.30\columnwidth}\raggedright
filter through cloth, then decant supernatant after 24 h settling\strut
\end{minipage}\tabularnewline
\begin{minipage}[t]{0.30\columnwidth}\raggedright
𐑫\strut
\end{minipage} & \begin{minipage}[t]{0.30\columnwidth}\raggedright
full chiral separation\strut
\end{minipage} & \begin{minipage}[t]{0.30\columnwidth}\raggedright
preparative column or liquid--liquid partition (four-step process)\strut
\end{minipage}\tabularnewline
\bottomrule
\end{longtable}

Ħ=𐑫 (full chiral separation) co-occurs with ⊙=𐑮 (near-critical endpoint,
\textless2\%) in all observed entries. This pairing is structurally
forced: near-critical endpoint monitoring requires
chiral-resolution-grade clarification to discriminate stereoisomeric
fractions at the sub-2\% level.

\hypertarget{ux262-coupling-fraction-combination}{%
\subsubsection{ɢ --- Coupling → Fraction
Combination}\label{ux262-coupling-fraction-combination}}

\begin{longtable}[]{@{}lll@{}}
\toprule
\begin{minipage}[b]{0.30\columnwidth}\raggedright
Shavian\strut
\end{minipage} & \begin{minipage}[b]{0.30\columnwidth}\raggedright
Mode\strut
\end{minipage} & \begin{minipage}[b]{0.30\columnwidth}\raggedright
Description\strut
\end{minipage}\tabularnewline
\midrule
\endhead
\begin{minipage}[t]{0.30\columnwidth}\raggedright
𐑝\strut
\end{minipage} & \begin{minipage}[t]{0.30\columnwidth}\raggedright
sequential\strut
\end{minipage} & \begin{minipage}[t]{0.30\columnwidth}\raggedright
add fractions one after another; evaluate each before combining\strut
\end{minipage}\tabularnewline
\begin{minipage}[t]{0.30\columnwidth}\raggedright
𐑜\strut
\end{minipage} & \begin{minipage}[t]{0.30\columnwidth}\raggedright
paired\strut
\end{minipage} & \begin{minipage}[t]{0.30\columnwidth}\raggedright
combine in pairs; evaluate paired yield\strut
\end{minipage}\tabularnewline
\begin{minipage}[t]{0.30\columnwidth}\raggedright
𐑠\strut
\end{minipage} & \begin{minipage}[t]{0.30\columnwidth}\raggedright
parallel\strut
\end{minipage} & \begin{minipage}[t]{0.30\columnwidth}\raggedright
combine all fractions simultaneously in a single vessel\strut
\end{minipage}\tabularnewline
\begin{minipage}[t]{0.30\columnwidth}\raggedright
𐑵\strut
\end{minipage} & \begin{minipage}[t]{0.30\columnwidth}\raggedright
broadcast\strut
\end{minipage} & \begin{minipage}[t]{0.30\columnwidth}\raggedright
broadcast combined fraction to multiple preparation vessels\strut
\end{minipage}\tabularnewline
\bottomrule
\end{longtable}

\hypertarget{criticality-endpoint-criterion}{%
\subsubsection{⊙ --- Criticality → Endpoint
Criterion}\label{criticality-endpoint-criterion}}

Criticality encodes when to stop the extraction. The Frobenius fixed
point at ⊙=⊙ is the structural self-referential endpoint: the extraction
has reached the point where successive fractions reproduce the same
pattern.

\begin{longtable}[]{@{}lll@{}}
\toprule
Shavian & Endpoint & Criterion\tabularnewline
\midrule
\endhead
𐑢 & sub-critical & stop before saturation; 70--80\% extraction
efficiency\tabularnewline
⊙ & at criticality & Frobenius fixed point; successive fractions differ
\textless{} 5\%\tabularnewline
𐑮 & near-critical & continue past threshold; successive fractions differ
\textless{} 2\%\tabularnewline
𐑻 & super-critical & drive to completion; \textless{} 1\% residual in
marc\tabularnewline
𐑣 & hyper-critical & exhaustive extraction; marc assayed for residual
content\tabularnewline
\bottomrule
\end{longtable}

\hypertarget{uxf0-dimensionality-and-ux159-recognition}{%
\subsubsection{Ð --- Dimensionality and Ř ---
Recognition}\label{uxf0-dimensionality-and-ux159-recognition}}

These two primitives are structural address parameters rather than
direct pharmaceutical parameters. Ð encodes the registration depth
(point, linear, planar, volumetric, etc.) and Ř encodes the
pattern-completion class of the entry. Both inform the session routing
--- particularly Gate 1 distance computation --- but do not directly
annotate recipe steps.

\begin{center}\rule{0.5\linewidth}{0.5pt}\end{center}

\hypertarget{the-recipe-opcodes}{%
\subsection{The Recipe Opcodes}\label{the-recipe-opcodes}}

The VMS recipe section uses a small, closed set of pharmaceutical
operation names. Each opcode names an abstract operation. The IG tuple
supplies its parameters.

\begin{longtable}[]{@{}lll@{}}
\toprule
\begin{minipage}[b]{0.30\columnwidth}\raggedright
Opcode\strut
\end{minipage} & \begin{minipage}[b]{0.30\columnwidth}\raggedright
Operation\strut
\end{minipage} & \begin{minipage}[b]{0.30\columnwidth}\raggedright
Primitives Consulted\strut
\end{minipage}\tabularnewline
\midrule
\endhead
\begin{minipage}[t]{0.30\columnwidth}\raggedright
Accipe\strut
\end{minipage} & \begin{minipage}[t]{0.30\columnwidth}\raggedright
Receive / take the plant material\strut
\end{minipage} & \begin{minipage}[t]{0.30\columnwidth}\raggedright
Þ (material specification)\strut
\end{minipage}\tabularnewline
\begin{minipage}[t]{0.30\columnwidth}\raggedright
Divide\strut
\end{minipage} & \begin{minipage}[t]{0.30\columnwidth}\raggedright
Comminute the material\strut
\end{minipage} & \begin{minipage}[t]{0.30\columnwidth}\raggedright
Γ (mesh/size)\strut
\end{minipage}\tabularnewline
\begin{minipage}[t]{0.30\columnwidth}\raggedright
Tere\strut
\end{minipage} & \begin{minipage}[t]{0.30\columnwidth}\raggedright
Triturate / grind\strut
\end{minipage} & \begin{minipage}[t]{0.30\columnwidth}\raggedright
Γ (mesh/size)\strut
\end{minipage}\tabularnewline
\begin{minipage}[t]{0.30\columnwidth}\raggedright
Extrahe\strut
\end{minipage} & \begin{minipage}[t]{0.30\columnwidth}\raggedright
Extract\strut
\end{minipage} & \begin{minipage}[t]{0.30\columnwidth}\raggedright
Φ (solvent), Σ (ratio), Ω (cycles), ⊙ (endpoint)\strut
\end{minipage}\tabularnewline
\begin{minipage}[t]{0.30\columnwidth}\raggedright
Calefac\strut
\end{minipage} & \begin{minipage}[t]{0.30\columnwidth}\raggedright
Heat\strut
\end{minipage} & \begin{minipage}[t]{0.30\columnwidth}\raggedright
Ç (process/temperature); ⚠ see cold-process constraint\strut
\end{minipage}\tabularnewline
\begin{minipage}[t]{0.30\columnwidth}\raggedright
Commisce\strut
\end{minipage} & \begin{minipage}[t]{0.30\columnwidth}\raggedright
Mix / combine\strut
\end{minipage} & \begin{minipage}[t]{0.30\columnwidth}\raggedright
ɢ (combination mode)\strut
\end{minipage}\tabularnewline
\begin{minipage}[t]{0.30\columnwidth}\raggedright
Colare\strut
\end{minipage} & \begin{minipage}[t]{0.30\columnwidth}\raggedright
Filter / clarify\strut
\end{minipage} & \begin{minipage}[t]{0.30\columnwidth}\raggedright
Ħ (clarification), + Φ/Σ/Ω/⊙ context\strut
\end{minipage}\tabularnewline
\begin{minipage}[t]{0.30\columnwidth}\raggedright
Compone\strut
\end{minipage} & \begin{minipage}[t]{0.30\columnwidth}\raggedright
Compose / endpoint\strut
\end{minipage} & \begin{minipage}[t]{0.30\columnwidth}\raggedright
ɢ (combination), ⊙ (endpoint criterion)\strut
\end{minipage}\tabularnewline
\begin{minipage}[t]{0.30\columnwidth}\raggedright
Transmuta\strut
\end{minipage} & \begin{minipage}[t]{0.30\columnwidth}\raggedright
Volatile transformation\strut
\end{minipage} & \begin{minipage}[t]{0.30\columnwidth}\raggedright
Ç (process); ⚠ cold-process warning if applicable\strut
\end{minipage}\tabularnewline
\begin{minipage}[t]{0.30\columnwidth}\raggedright
Applica\strut
\end{minipage} & \begin{minipage}[t]{0.30\columnwidth}\raggedright
Apply / administer\strut
\end{minipage} & \begin{minipage}[t]{0.30\columnwidth}\raggedright
ƒ (concentration and dosing form)\strut
\end{minipage}\tabularnewline
\begin{minipage}[t]{0.30\columnwidth}\raggedright
Administra\strut
\end{minipage} & \begin{minipage}[t]{0.30\columnwidth}\raggedright
Administer\strut
\end{minipage} & \begin{minipage}[t]{0.30\columnwidth}\raggedright
ƒ (concentration and dosing form)\strut
\end{minipage}\tabularnewline
\bottomrule
\end{longtable}

\hypertarget{calefac-under-cold-process-constraint}{%
\subsubsection{Calefac Under Cold-Process
Constraint}\label{calefac-under-cold-process-constraint}}

When a plant's Ç value is cold maceration (𐑤), the Calefac opcode cannot
refer to the primary extract --- it is structurally incompatible. The
step is annotated:

\begin{verbatim}
Process (adjunct): heat excipient / base as needed
⚠ Primary extract: cold maceration — do not heat
\end{verbatim}

The heat is applied to the carrier medium, the binding agent, or the
secondary component that receives the cold extract. This is not an
interpretation; it is the only structurally consistent reading when Ç=𐑤.

\hypertarget{transmuta-under-cold-process-constraint}{%
\subsubsection{Transmuta Under Cold-Process
Constraint}\label{transmuta-under-cold-process-constraint}}

Transmuta (volatile phase transformation) is similarly annotated with a
warning for cold-process plants, since the volatile separation step
involves elevated temperatures that would damage a frozen-order extract.
The cold extract should be introduced after the volatile fraction has
been stabilized.

\begin{center}\rule{0.5\linewidth}{0.5pt}\end{center}

\hypertarget{the-session-engine-six-gates}{%
\subsection{The Session Engine --- Six
Gates}\label{the-session-engine-six-gates}}

The session engine is the linker that connects a plant's structural
identity to a VMS recipe and then elaborates that recipe with the
plant's protocol parameters.

\hypertarget{init-cosmological-initialization}{%
\subsubsection{INIT --- Cosmological
Initialization}\label{init-cosmological-initialization}}

Source: f68 (cosmological foldout, the largest VMS folio).

The foldout establishes two structural invariants that persist across
all subsequent gates:

\begin{itemize}
\item
  \textbf{Φ (Parity)}: The solvent class is conferred physically --- the
  practitioner selects the appropriate solvent medium before beginning.
  This corresponds to the foldout's role as the ``pre-session''
  initialization: it cannot be derived from the recipe, it must be known
  in advance.
\item
  \textbf{Ħ (Chirality)}: The clarification protocol is similarly fixed
  at initialization. The practitioner commits to a separation strategy
  before any extraction begins.
\end{itemize}

Both invariants are read from the plant's IG tuple. The foldout is the
structural site where the tuple is ``loaded'' into the session.

\hypertarget{gate-1-pharmaceutical-address}{%
\subsubsection{GATE 1 --- Pharmaceutical
Address}\label{gate-1-pharmaceutical-address}}

Source: Pharmacy section (f99--f102).

Selection criteria: \textbf{potency class} and \textbf{pars plantae}
(plant part). These are the only two selection keys. Applicatio (route
of administration) is a kinetics descriptor carried in Ç --- it is NOT a
pharmacy selection criterion.

\begin{itemize}
\tightlist
\item
  Summa potency + folium/flos → \textbf{f11r/p6} (leaf/flower summa;
  extractio → mixtura; n\_ops=11; fixatio=True)
\item
  Summa potency + radix → \textbf{f39r/p3} (root/rhizome summa;
  trituratio + extractio → mixtura; n\_ops=12; fixatio=False)
\end{itemize}

Gate passes when a unique pharmacy entry matches the potency and part.
Failure here means the plant is not addressable by the standard pharmacy
catalog --- a structural incompatibility, not a session error.

\hypertarget{gate-2-balneological-heap}{%
\subsubsection{GATE 2 --- Balneological
Heap}\label{gate-2-balneological-heap}}

Source: Balneological section (f75--f84); 20 folios indexed f75r, f75v,
\ldots, f84r, f84v.

Selection:
\texttt{heap\_folio\ =\ folios{[}pharmacy\_entry.folio\_number\ \%\ 20{]}}

The balneological section represents the ``containment vessel'' --- the
physical and thermal environment that receives the pharmaceutical
address. The heap folio must satisfy:

\begin{itemize}
\item
  \textbf{FSPLIT ≥ n\_ops}: The vessel's split capacity (number of
  distinct process branches) must be at least equal to the operation
  count of the pharmacy address. This verifies the vessel can hold the
  full instruction depth.
\item
  \textbf{FFUSE/FSPLIT ≥ 0.60} (non-volatile preparations only): For
  preparations that do not involve a volatile transformation, the vessel
  must be predominantly closed (fused). Oil and decoction entries carry
  this constraint; cold maceration entries with a volatile step do not.
\end{itemize}

Gate passes when both conditions are satisfied. This gate has never
failed in observed runs --- it appears to function as a structural
invariant rather than a contingent filter.

\hypertarget{gate-3-astronomical-winding}{%
\subsubsection{GATE 3 --- Astronomical
Winding}\label{gate-3-astronomical-winding}}

Source: Astronomical section, f69; three independent transcription
sources (H=Hermetic, F=Folkwang, U=University).

Selection: \texttt{paragraph\ =\ folio\_number\ \%\ 4} from f69.

The astronomical section is the winding register --- it provides
external, non-botanical authority for the cycle count and temporal
structure of the extraction. The three transcription sources are
independent attestations. The gate runs EVALT (evaluation of
transcription alignment) across all three:

\begin{itemize}
\tightlist
\item
  Where all three sources agree: the structural value is verified.
\item
  Where sources disagree: the disagreement position IS the structural
  signal. The transcribers independently marked the same decision point
  (VINIT/FFUSE ambiguity) without knowing it. Majority vote determines
  the gate value.
\end{itemize}

Gate passes when EVALT majority yields a consistent winding class. The Ω
value derived here confirms or overrides the Ω from the IG tuple. In
practice, they agree --- the catalog and the astronomical register are
coherent.

\hypertarget{output-recipe-selection}{%
\subsubsection{OUTPUT --- Recipe
Selection}\label{output-recipe-selection}}

Source: Recipe section (f103r+).

The recipe folios closest to the Gate 1 pharmacy address are ranked by
folio proximity. The top three recipes are returned. In all nine
inaugural runs:

\begin{itemize}
\tightlist
\item
  \textbf{f103r/p2} --- 15 steps; extract-dominant path; opens with
  Extrahe
\item
  \textbf{f103r/p9} --- 12 steps; ingredient-led path; opens with Accipe
\item
  \textbf{f103r/p49} --- 12 steps; volatile-phase path; opens with
  Transmuta
\end{itemize}

These three represent three canonical execution paths through the same
pharmaceutical operation class. They are not alternatives --- they are
sequential phases of a complete preparation: bulk extraction (p2),
constituent assembly (p9), volatile transformation and final form (p49).

\hypertarget{elaboration-protocol-annotation}{%
\subsubsection{ELABORATION --- Protocol
Annotation}\label{elaboration-protocol-annotation}}

Once the recipe is selected, every opcode in every step is annotated
with the plant's protocol parameters derived from the IG tuple. The
annotation is deterministic: given the tuple, every parameter of every
step is fully specified. There are no free variables remaining after
elaboration.

\begin{center}\rule{0.5\linewidth}{0.5pt}\end{center}

\hypertarget{operation-end-to-end-run}{%
\subsection{Operation: End-to-End Run}\label{operation-end-to-end-run}}

A complete Voynich Phytoglyphica session for a botanical entry proceeds
as follows.

\textbf{Step 1: Catalog lookup.} The plant name is resolved against the
IG catalog (voynich\_pharmacy.json, 1491 entries). The 12-primitive
tuple is retrieved. Distances to all six VMS section tuples are computed
in 12-dimensional primitive space.

\textbf{Step 2: Protocol derivation.} The tuple is passed to the
elaborator. All ten pharmaceutical primitives (Þ, Φ, Ç, Σ, Γ, Ω, ƒ, Ħ,
ɢ, ⊙) are resolved against the lookup tables above. The protocol dict is
built. This happens before any gate is evaluated --- the full
preparation protocol is known from the tuple alone.

\textbf{Step 3: GATE 1.} Potency class is determined from the plant's
distance to the astronomical section tuple. If d(plant, astronomical) =
0, potency = summa. Part (pars plantae) is retrieved from the catalog.
The pharmacy entry is selected.

\textbf{Step 4: GATE 2.} The balneological heap folio is selected by
index. FSPLIT and FFUSE conditions are evaluated. Gate passes or fails.

\textbf{Step 5: GATE 3.} The astronomical paragraph is selected from
f69. EVALT majority across H/F/U sources is computed. Gate passes or
fails.

\textbf{Step 6: Recipe output.} If all three gates pass, up to three
recipe folios are selected by proximity. Steps are extracted from the
recipe database.

\textbf{Step 7: Elaborated output.} Every recipe step is printed with
its IG-derived annotations. The full protocol header (tuple, all
parameters, their descriptions) is printed before the recipe. The
practitioner now has a complete, quantified pharmaceutical instruction
set.

\begin{center}\rule{0.5\linewidth}{0.5pt}\end{center}

\hypertarget{implications}{%
\subsection{Implications}\label{implications}}

\hypertarget{the-vms-is-not-a-failed-cipher}{%
\subsubsection{The VMS Is Not a Failed
Cipher}\label{the-vms-is-not-a-failed-cipher}}

The VMS resisted decipherment because the approach of treating it as a
natural language text was structurally incorrect. The manuscript does
not contain its own parameters. It was designed to be read alongside a
structural grammar that the practitioners already possessed. The opcodes
(Extrahe, Calefac, Colare) are unambiguous pharmaceutical Latin; they
were never the problem. The parameters were the problem, and they were
never in the manuscript.

\hypertarget{session-invariance-protocol-variance}{%
\subsubsection{Session Invariance / Protocol
Variance}\label{session-invariance-protocol-variance}}

The session engine produces identical gate behavior across all entries
tested. All nine inaugural plants pass all three gates and receive the
same three recipe folios. The differentiation between preparations lives
entirely in the IG tuple. This is architecturally correct: the VMS
separates \emph{what class of operation} to perform (session routing)
from \emph{how to perform it} (tuple elaboration). A practitioner with
the wrong tuple would produce the right sequence of operations with the
wrong parameters --- pharmacologically inert or hazardous.

\hypertarget{the-wormwood-demonstration}{%
\subsubsection{The Wormwood
Demonstration}\label{the-wormwood-demonstration}}

Artemisia absinthium appears twice in the inaugural set with ten of
twelve identical primitive values. The single difference is ƒ: Fidelity
𐑱 (1×, proportional) vs 𐑐 (3×, standardized concentrate). The same
plant, the same extraction sequence, but one preparation is a
proportional tincture and the other is a reduced, standardized,
marker-verified concentrate. The grammar distinguishes what morphology
cannot. This is the key empirical demonstration: the IG provides
pharmaceutical discrimination at a resolution beyond botanical taxonomy.

\hypertarget{the-cold-process-constraint-as-structural-proof}{%
\subsubsection{The Cold-Process Constraint as Structural
Proof}\label{the-cold-process-constraint-as-structural-proof}}

The Calefac-Ç interaction is not a heuristic. It is a structural
entailment: if Ç=𐑤, then any Calefac in the recipe must be assigned to a
secondary component, because the tuple and the opcode are structurally
inconsistent if assigned to the same object. This forces a reparse of
the recipe step --- one that assigns Calefac to the excipient and not
the extract. The grammar performs this reparse automatically. A reader
without the grammar would see a heating instruction applied to a plant
that cannot be heated and have no structural basis for the correct
interpretation.

\hypertarget{the-vms-as-universal-engine}{%
\subsubsection{The VMS as Universal
Engine}\label{the-vms-as-universal-engine}}

The session architecture is plant-agnostic at the gate level. Any IG
catalog entry with a pharmaceutical tuple can be processed. The engine
does not know about wormwood specifically; it knows about topology,
kinetics, stoichiometry, winding. The VMS is not a recipe book for
specific plants --- it is a universal pharmaceutical engine whose
parameters are supplied by the structural grammar of whatever is being
prepared. This is why it resisted decipherment as a specific-language
text: it is not specific to any language, any plant, or any culture. It
is a grammar-parameterized computation.

The Imscribing Grammar is the parameter set. The Voynich Manuscript is
the program.

\begin{center}\rule{0.5\linewidth}{0.5pt}\end{center}

Voynich Phytoglyphica Engine Specification Source:
voynich\_phytoglyphica package; voynich\_engine session engine IG
catalog (voynich\_pharmacy.json); LSI\_ivtff\_0d.txt transcription
